\documentclass{article}
\usepackage{/globals/anki}
\ankiprefix{staatsorgane}
\ankitopic{Staatsorgane}
\ankishow
\begin{document}
Eine im Grundgesetz verankerte rechtlich geschaffene Einrichtung die für einen Rechtsträger handelt, mit bestimmten Rechten und Pflichten, ist ein \emph{Staatsorgan}, auch \emph{Verfasungsorgan} genannt.
\ankinote{def}{Defitition}{rechtlich geschaffene Einrichtung handelt für Rechtsträger; mit Rechten und Pflichten}

Dieses hat in der Regel eine rechtliche \emph{Hauptfunktion}; die Gesetze die für das Organ wichtig sind sind dessen \emph{Rechtsquellen}.
\ankinote{hauptfunktion}{Hauptfunktion}{Gesetzliche zentrale Aufgabe}
\ankinote{relevante-gesetze}{Relevante Gesetze}{Rechtsquellen}

Für uns relevant sind (ständig) der Bundestag, Bundesrat, der Bundespräsident, die Bundesregierung, das Bundesverfassungsgericht und (nicht ständig) der Gemeinsame Ausschuss, der Vermittlungsausschuss und die Bundesversammlung.
\ankinote{gesamt}{Gesamt}{8}
\ankinote{ständig}{Ständig}{5}

Alle Organe handeln natürlich nach der Gewaltenteilung; sowohl horizontal (Teilung in Legislative, Judikative und Exekutive)
\ankinote{gewalt-horiontal}{Gewaltenteilung Horizontal}{Drei Gewalten}
\ankinote{gewalt-vertikal}{Gewaltenteilung Vertikal}{Bund, Länder}
\end{document}